\documentclass[11pt,a4paper]{article}
\usepackage[left=24mm,right=24mm,top=22mm,bottom=20mm,includefoot]{geometry}
\usepackage{amsmath,amssymb,booktabs}
\usepackage{iftex}
\ifPDFTeX
  \usepackage{CJKutf8}
  % CJKutf8 supplies its own Unicode map for each Chinese subfont.
  \makeatletter
  \AddToHook{cmd/CJK@addcmap/before}{\pdfnobuiltintounicode\font@name}
  \makeatother
\else
  \usepackage[no-math]{fontspec}
  \IfFontExistsTF{Noto Serif CJK SC}{
    \setmainfont{Noto Serif CJK SC}
    \setsansfont{Noto Sans CJK SC}
  }{
    \setmainfont{FandolSong-Regular.otf}[BoldFont=FandolSong-Bold.otf]
    \setsansfont{FandolHei-Regular.otf}[BoldFont=FandolHei-Bold.otf]
  }
  \ifXeTeX
    \XeTeXlinebreaklocale "zh"
    \XeTeXlinebreakskip=0pt plus 0.5pt
  \fi
\fi
\pagestyle{plain}
\raggedbottom
\linespread{1.13}
\setlength{\parindent}{2em}
\setlength{\parskip}{3pt plus 1pt}
\setlength{\jot}{3pt}
\DeclareMathSizes{10.95}{10.95}{9}{7}
% Use a single stroke for Chinese headings when pdfLaTeX has no bold CJK font.
\newcommand{\headingweight}{\ifPDFTeX\mdseries\else\bfseries\fi}
\makeatletter
\renewcommand{\section}{\@startsection{section}{1}{0pt}
  {2.0ex plus .3ex minus .2ex}{.9ex}{\normalfont\large\headingweight}}
\makeatother

\newcommand{\indep}{\mathrel{\perp\!\!\!\perp}}
\begin{document}
\ifPDFTeX\begin{CJK*}{UTF8}{gbsn}\fi
\begin{center}
{\LARGE\headingweight 两个暴露的有序结局 Shapley 归因}\\[5pt]
{\large 单步增量与 Markov 条件下的识别}
\end{center}

\section{要分配的责任是什么}
设 $X_1,X_2\in\{0,1\}$，结局 $Y\in\{0,1,2\}$，分别表示小、中、大。
每升一级计 $1$ 分；这里分配的是等级分数。四个潜在结果的含义为
\begin{center}
\renewcommand{\arraystretch}{1.17}
\begin{tabular}{@{}c c l@{}}
\toprule
暴露组合 $(X_1,X_2)$ & 潜在结果 & 含义\\
\midrule
$(0,0)$ & $Y_{00}$ & 两个暴露均不存在\\
$(0,1)$ & $Y_{01}$ & 只有暴露 $X_2$\\
$(1,0)$ & $Y_{10}$ & 只有暴露 $X_1$\\
$(1,1)$ & $Y_{11}$ & 两个暴露均存在\\
\bottomrule
\end{tabular}
\end{center}
个体类型按 $T=(Y_{00},Y_{01},Y_{10},Y_{11})$ 排列。定义
\[
\mathcal E_1=\{X_1=X_2=1,Y=1\},\qquad
\mathcal E_2=\{X_1=X_2=1,Y=2\}.
\]
由一致性，在 $\mathcal E_2$ 中 $Y_{11}=2$。这 $2$ 分由无暴露时已有的 $Y_{00}$ 分，
以及两个暴露带来的 $2-Y_{00}$ 分组成。

\medskip
\noindent{\headingweight 给定类型的责任。}
Shapley 规则对两种加入顺序各取一半权重。
\begin{center}
\renewcommand{\arraystretch}{1.22}
\begin{tabular}{@{}c c c c@{}}
\toprule
加入顺序 & 潜在结果路径 & $X_1$ 的增量 & $X_2$ 的增量\\
\midrule
先 $X_1$，后 $X_2$ & $Y_{00}\to Y_{10}\to Y_{11}$
 & $Y_{10}-Y_{00}$ & $Y_{11}-Y_{10}$\\
先 $X_2$，后 $X_1$ & $Y_{00}\to Y_{01}\to Y_{11}$
 & $Y_{11}-Y_{01}$ & $Y_{01}-Y_{00}$\\
\bottomrule
\end{tabular}
\end{center}
用编号 $0$ 表示无暴露背景。对给定 $T=t$，将各潜在结果取为 $t$ 中对应的值，定义
\begin{equation}
\begin{aligned}
R(1\mid t,\mathcal E_2)&=\tfrac12\{(Y_{10}-Y_{00})+(Y_{11}-Y_{01})\},\\
R(2\mid t,\mathcal E_2)&=\tfrac12\{(Y_{01}-Y_{00})+(Y_{11}-Y_{10})\},\\
R(0\mid t,\mathcal E_2)&=Y_{00}.
\end{aligned}
\label{eq:individual}
\end{equation}
三项合计为 $Y_{11}=2$。在 $\mathcal E_1$ 中使用同一规则，三项合计为 $1$。

\medskip
\noindent{\headingweight 病例中的平均责任。}
例如，对 $j=0,1,2$，
\[
\mathcal R_j(\mathcal E_2)
=\sum_{t:\,t_{11}=2}R(j\mid t,\mathcal E_2)\,
  \mathrm{pr}(T=t\mid\mathcal E_2).
\]
识别的目标是这个条件平均值。以下假设一致性成立，联合暴露 $(X_1,X_2)$ 与 $T$ 独立，
四个暴露组均有正概率，且所讨论的病例事件发生概率为正。
于是条件事件 $\mathcal E_2$ 可替换为 $Y_{11}=2$；$\mathcal E_1$ 同理。

\clearpage
\section{需要增加哪两个条件}
\noindent{\headingweight 条件一：每加入一个暴露，结局最多升一级。}
要求同一个体的四个潜在结果满足
\begin{equation}
\begin{aligned}
0\le Y_{10}-Y_{00}\le1,\qquad&0\le Y_{11}-Y_{10}\le1,\\
0\le Y_{01}-Y_{00}\le1,\qquad&0\le Y_{11}-Y_{01}\le1.
\end{aligned}
\label{eq:unit}
\end{equation}
这允许两个暴露合起来升两级。例如 $T=0112$ 对应两条路径均为 $0\to1\to2$。

\medskip
\noindent{\headingweight 条件二：给定中间等级，最终等级与起始等级条件独立。}
对两种加入顺序，分别写为
\begin{equation}
\begin{aligned}
Y_{11}\indep Y_{00}\mid Y_{10},&\qquad
    \text{对应 }Y_{00}\to Y_{10}\to Y_{11},\\
Y_{11}\indep Y_{00}\mid Y_{01},&\qquad
    \text{对应 }Y_{00}\to Y_{01}\to Y_{11}.
\end{aligned}
\label{eq:markov}
\end{equation}
第一条的含义是：在 $Y_{10}=b$ 的人中，知道 $Y_{00}=a$ 不再改变 $Y_{11}$ 的条件分布，
即在条件事件概率为正时，
\[
\mathrm{pr}(Y_{11}=c\mid Y_{10}=b,Y_{00}=a)
=\mathrm{pr}(Y_{11}=c\mid Y_{10}=b).
\]
第二条将 $Y_{10}$ 换成 $Y_{01}$。这些是跨潜在结果的额外假设，随机分配本身不蕴含它们。
路径表示计算 Shapley 值时的加入顺序，不表示实际随访时间。

两条路径均只有起始、中间和最终三个潜在结果，因此式~\eqref{eq:markov} 就是相应路径的 Markov 条件。

\section{第一步：恢复相邻两个潜在结果的联合分布}
以 $Y_{00}$ 与 $Y_{10}$ 为例。条件 $0\le Y_{10}-Y_{00}\le1$ 限制了二者的联合取值。
下表列出全部五种可能，其概率均由这两个潜在结果各自的分布确定：
\begin{center}
\renewcommand{\arraystretch}{1.35}
\setlength{\tabcolsep}{12pt}
\begin{tabular}{@{}c l@{}}
\toprule
$(Y_{00},Y_{10})$ 的取值 & 对应的联合概率\\
\midrule
$(0,0)$ & $\mathrm{pr}(Y_{10}=0)$\\
$(0,1)$ & $\mathrm{pr}(Y_{00}=0)-\mathrm{pr}(Y_{10}=0)$\\
$(1,1)$ & $1-\mathrm{pr}(Y_{10}=2)-\mathrm{pr}(Y_{00}=0)$\\
$(1,2)$ & $\mathrm{pr}(Y_{10}=2)-\mathrm{pr}(Y_{00}=2)$\\
$(2,2)$ & $\mathrm{pr}(Y_{00}=2)$\\
\bottomrule
\end{tabular}
\end{center}
例如，$Y_{00}=0$ 的人加入 $X_1$ 后只能停在 $0$ 或升至 $1$。
其中停在 $0$ 的概率就是 $\mathrm{pr}(Y_{10}=0)$，
因此从 $0$ 升至 $1$ 的概率为 $\mathrm{pr}(Y_{00}=0)-\mathrm{pr}(Y_{10}=0)$。
从 $1$ 升至 $2$ 的概率同理。

同样的计算分别用于 $(Y_{10},Y_{11})$、$(Y_{00},Y_{01})$ 和 $(Y_{01},Y_{11})$，
便恢复其余三对相邻联合分布。
四个潜在结果的边际概率由相应随机暴露组的结局比例识别。

\clearpage
\section{第二步：恢复无暴露与双暴露的联合分布}
第一步已经给出 $Y_{10}$ 与 $Y_{11}$、$Y_{01}$ 与 $Y_{11}$ 的联合分布。
还缺的是 $Y_{00}$ 与 $Y_{11}$ 的联合分布，因为两者相隔两步。
沿 $Y_{00}\to Y_{10}\to Y_{11}$ 使用第一条 Markov 条件。当 $\mathrm{pr}(Y_{10}=b)>0$ 时，
\begin{equation}
\begin{aligned}
&\mathrm{pr}(Y_{00}=a,Y_{10}=b,Y_{11}=c)\\
&\quad=\mathrm{pr}(Y_{00}=a,Y_{10}=b)\,
       \mathrm{pr}(Y_{11}=c\mid Y_{00}=a,Y_{10}=b)\\
&\quad=\mathrm{pr}(Y_{00}=a,Y_{10}=b)\,
       \mathrm{pr}(Y_{11}=c\mid Y_{10}=b)\\
&\quad=\frac{\mathrm{pr}(Y_{00}=a,Y_{10}=b)\,
             \mathrm{pr}(Y_{10}=b,Y_{11}=c)}
            {\mathrm{pr}(Y_{10}=b)}.
\end{aligned}
\label{eq:triple}
\end{equation}
中间那个等号正是 Markov 条件发挥作用的地方。
对 $b=0,1,2$ 求和，得到
\begin{equation}
\mathrm{pr}(Y_{00}=a,Y_{11}=c)
=\sum_{b:\,\mathrm{pr}(Y_{10}=b)>0}
\frac{\mathrm{pr}(Y_{00}=a,Y_{10}=b)\,
      \mathrm{pr}(Y_{10}=b,Y_{11}=c)}
     {\mathrm{pr}(Y_{10}=b)}.
\label{eq:endpoints}
\end{equation}
分母为零的层对联合概率贡献为零，故不进入求和。

\medskip
\noindent{\headingweight 对二级病例，最关键的是从 $0$ 到 $2$ 的概率。}
记
\[
\begin{aligned}
\xi_{00}&=\mathrm{pr}(Y_{00}=2),&\qquad \xi_{01}&=\mathrm{pr}(Y_{01}=2),\\
\xi_{10}&=\mathrm{pr}(Y_{10}=2),&\qquad \xi_{11}&=\mathrm{pr}(Y_{11}=2),
\end{aligned}
\]
以及 $\theta=\mathrm{pr}(Y_{00}=0,Y_{11}=2)$。

在条件~\eqref{eq:unit} 下，$Y_{00}=0,Y_{11}=2$ 强制两条路径的中间值均为 $1$，因此
\[
\{Y_{00}=0,Y_{11}=2\}=\{T=0112\}.
\]
于是式~\eqref{eq:endpoints} 在 $a=0,c=2$ 时只剩 $b=1$ 这一项：
\[
\theta=
\frac{\mathrm{pr}(Y_{00}=0,Y_{10}=1)\,
      \mathrm{pr}(Y_{10}=1,Y_{11}=2)}
     {\mathrm{pr}(Y_{10}=1)}.
\]
第一步的表给出
\[
\begin{aligned}
\mathrm{pr}(Y_{00}=0,Y_{10}=1)
 &=\mathrm{pr}(Y_{00}=0)-\mathrm{pr}(Y_{10}=0),\\
\mathrm{pr}(Y_{10}=1,Y_{11}=2)&=\xi_{11}-\xi_{10}.
\end{aligned}
\]
所以当 $\mathrm{pr}(Y_{10}=1)>0$ 时，
\begin{equation}
\theta=
\frac{\{\mathrm{pr}(Y_{00}=0)-\mathrm{pr}(Y_{10}=0)\}
      (\xi_{11}-\xi_{10})}
     {\mathrm{pr}(Y_{10}=1)}.
\label{eq:theta}
\end{equation}
若 $\mathrm{pr}(Y_{10}=1)=0$，则 $\theta=0$。
式~\eqref{eq:theta} 右侧全部是可由四个随机暴露组识别的边际概率。

\clearpage
\section{第三步：计算二级病例的责任}
在单步增量条件下，$\mathcal E_2$ 只允许以下六种类型。两项暴露责任仍按式~\eqref{eq:individual}
计算，并未改变原来的 Shapley 规则。
\begin{center}
\renewcommand{\arraystretch}{1.28}
\setlength{\tabcolsep}{14pt}
\begin{tabular}{@{}c c c c@{}}
\toprule
类型 $t$ & $R(1\mid t,\mathcal E_2)$ & $R(2\mid t,\mathcal E_2)$
 & $R(0\mid t,\mathcal E_2)$\\
\midrule
$0112$ & $1$ & $1$ & $0$\\
$1112$ & $1/2$ & $1/2$ & $1$\\
$1122$ & $1$ & $0$ & $1$\\
$1212$ & $0$ & $1$ & $1$\\
$1222$ & $1/2$ & $1/2$ & $1$\\
$2222$ & $0$ & $0$ & $2$\\
\bottomrule
\end{tabular}
\end{center}

为了求这些行的加权平均，只需要三个条件均值。
当 $Y_{11}=2$ 时，由单步增量限制，$Y_{10}$ 只能是 $1$ 或 $2$。
其中取 $2$ 的条件概率为 $\xi_{10}/\xi_{11}$，所以
\[
\mathrm E(Y_{10}\mid Y_{11}=2)=1+\frac{\xi_{10}}{\xi_{11}},
\qquad
\mathrm E(Y_{01}\mid Y_{11}=2)=1+\frac{\xi_{01}}{\xi_{11}}.
\]
背景 $Y_{00}$ 在二级病例中的条件分布为
\[
\begin{aligned}
\mathrm{pr}(Y_{00}=0\mid Y_{11}=2)&=\frac{\theta}{\xi_{11}},\\
\mathrm{pr}(Y_{00}=1\mid Y_{11}=2)&=\frac{\xi_{11}-\theta-\xi_{00}}{\xi_{11}},\\
\mathrm{pr}(Y_{00}=2\mid Y_{11}=2)&=\frac{\xi_{00}}{\xi_{11}}.
\end{aligned}
\]
因此 $\mathrm E(Y_{00}\mid Y_{11}=2)=(\xi_{11}+\xi_{00}-\theta)/\xi_{11}$。

\medskip
\noindent{\headingweight 代入原始 Shapley 公式。}
例如，
\[
\mathcal R_1(\mathcal E_2)
=\tfrac12\{\mathrm E(Y_{10}\mid Y_{11}=2)-\mathrm E(Y_{00}\mid Y_{11}=2)
             +2-\mathrm E(Y_{01}\mid Y_{11}=2)\}.
\]
将三个条件均值代入并化简，得到
\begin{equation}
\begin{aligned}
\mathcal R_1(\mathcal E_2)
 &=\frac{\xi_{11}-\xi_{01}+\xi_{10}-\xi_{00}+\theta}{2\xi_{11}},\\[3pt]
\mathcal R_2(\mathcal E_2)
 &=\frac{\xi_{11}-\xi_{10}+\xi_{01}-\xi_{00}+\theta}{2\xi_{11}},\\[3pt]
\mathcal R_{0}(\mathcal E_2)
 &=\frac{\xi_{11}+\xi_{00}-\theta}{\xi_{11}}.
\end{aligned}
\label{eq:result}
\end{equation}
这三项均可点识别，合计为 $2$。其中四个 $\xi$ 概率直接来自各组的二级结局比例，
$\theta$ 由式~\eqref{eq:theta} 识别。证明不要求先识别每一种类型的条件概率。

\clearpage
\section{一个从观测分布算到责任的例子}
假设四个随机暴露组的结局分布如下。
\begin{center}
\renewcommand{\arraystretch}{1.20}
\setlength{\tabcolsep}{17pt}
\begin{tabular}{@{}c c c c@{}}
\toprule
暴露组合 & $Y=0$ & $Y=1$ & $Y=2$\\
\midrule
$(0,0)$ & $1/2$ & $1/2$ & $0$\\
$(1,0)$ & $0$ & $1$ & $0$\\
$(0,1)$ & $0$ & $1$ & $0$\\
$(1,1)$ & $0$ & $1/2$ & $1/2$\\
\bottomrule
\end{tabular}
\end{center}
于是 $\xi_{00}=\xi_{10}=\xi_{01}=0$，$\xi_{11}=1/2$。
因为 $Y_{10}=Y_{01}=1$，Markov 条件在本例中要求 $Y_{00}$ 与 $Y_{11}$ 独立。
按式~\eqref{eq:theta}，
\[
\theta=\frac{(1/2-0)(1/2-0)}{1}=\frac14.
\]
代入式~\eqref{eq:result}，
\[
\begin{aligned}
\mathcal R_1(\mathcal E_2)&=\frac{1/2+1/4}{2(1/2)}=\frac34,\\
\mathcal R_2(\mathcal E_2)&=\frac{1/2+1/4}{2(1/2)}=\frac34,\\
\mathcal R_{0}(\mathcal E_2)&=\frac{1/2-1/4}{1/2}=\frac12.
\end{aligned}
\]
即二级病例的 $2$ 分中，两个暴露各贡献 $0.75$ 分，背景贡献 $0.5$ 分。

本例中，四种类型 $0111,0112,1111,1112$ 各占总体的 $1/4$。
因此二级病例仅包含以下两种类型，直接加权得到相同结果。
\begin{center}
\renewcommand{\arraystretch}{1.23}
\setlength{\tabcolsep}{13pt}
\begin{tabular}{@{}c c c c c@{}}
\toprule
类型 $t$ & $\mathrm{pr}(T=t\mid\mathcal E_2)$ & $X_1$ 责任 & $X_2$ 责任 & 背景责任\\
\midrule
$0112$ & $1/2$ & $1$ & $1$ & $0$\\
$1112$ & $1/2$ & $1/2$ & $1/2$ & $1$\\
\midrule
加权平均 & & $3/4$ & $3/4$ & $1/2$\\
\bottomrule
\end{tabular}
\end{center}

\medskip
\noindent{\headingweight 观测到一级时。}
将式~\eqref{eq:endpoints} 中的 $c$ 取为 $1$，结合已识别的相邻联合分布，
即可得到 $\mathrm E(Y_{00}\mid Y_{11}=1)$、$\mathrm E(Y_{10}\mid Y_{11}=1)$
和 $\mathrm E(Y_{01}\mid Y_{11}=1)$。代入同一 Shapley 公式即可识别 $\mathcal E_1$ 下的责任。
在本例中，一级病例由 $0111,1111$ 各占一半组成，
故两个暴露各贡献 $1/4$ 分，背景贡献 $1/2$ 分，合计为 $1$。

\medskip
\noindent{\headingweight 两条路径的使用范围。}
对 $p=2$，四条单步增量限制加上第一条 Markov 条件，已经足以完成上述识别。
若同时要求第二条，则将式~\eqref{eq:endpoints} 中的 $Y_{10}$ 全部换成 $Y_{01}$，
对每个 $a,c$ 必须得到同一个端点联合概率。
这些一致性等式是模型的必要限制；本结果以两条条件在同一个联合分布中共同成立为前提。

\ifPDFTeX\end{CJK*}\fi
\end{document}
